\begin{pastexam}[label={px:ust-115-ee-ds-50}]{臺聯大電機類 115 學年度 第 50 題}
（單選；2 分）

The following pseudo code is a simple algorithm to generate a permutation of n items. Please fill in
the BLANK area in Line 6 with a correct formula such that \texttt{i <= j < n}
\begin{minted}[linenos, numbersep=1.5mm, xleftmargin=2em, fontsize=\small]{text}
function Uniform(m)
    return a random integer between 0 and m-1;

function Permutation(A,n)
    for i = 0 to n-2
        j=i+Uniform(_____________)
        SWAP(A[i],A[j])
\end{minted}

(A)~\texttt{n} \qquad (B)~\texttt{n-i} \qquad (C)~\texttt{i} \qquad (D)~\texttt{n-1} \qquad (E)~\texttt{n/2}
\end{pastexam}

\begin{pastsoln}
答案是 (B)。

要 $i \le j < n$，也就是 $j - i$ 落在 $0, 1, \ldots, n - 1 - i$，共 $n - i$ 種值。
\proc{Uniform}$(m)$ 回傳 $0..m-1$，所以 $m = n - i$，$j = i + \proc{Uniform}(n - i)$ 在
$i..n-1$ 裡均勻分布。

這就是 Fisher--Yates 洗牌：第 $i$ 輪從還沒定位的 $A[i..n-1]$ 裡均勻挑一個放到第 $i$ 格。 第 $0$ 輪
有 $n$ 種選擇、第 $1$ 輪 $n - 1$ 種、……、第 $n - 2$ 輪 $2$ 種，每一串選擇對應一個不同的排列，共
$n!$ 串、機率各是 $1/n!$，所以每個排列出現的機率相同（機率的基本概念見資結書第~A.5~節）。 最後一格
不必再挑，所以迴圈只到 $n - 2$。

逐項檢查：
\begin{itemize}
\item (A) \texttt{n}：$j$ 可到 $i + n - 1$，$i \ge 1$ 時超出陣列。 錯。
\item (B) \texttt{n-i}：對。
\item (C) \texttt{i}：$i = 0$ 時 \proc{Uniform}$(0)$ 沒有可回傳的值；之後 $j$ 在 $i..2i-1$，
$i > n/2$ 時會超出 $n - 1$，而且挑不到 $2i$ 以後的位置。 錯。
\item (D) \texttt{n-1}：$j$ 可到 $i + n - 2$，$i \ge 2$ 時超出陣列。 錯。
\item (E) \texttt{n/2}：$i > n/2$ 時 $j$ 可能超出陣列；而且 $i$ 小時挑不到後半段，排列不均勻。 錯。
\end{itemize}
\end{pastsoln}
